\documentclass[11pt]{article}
\usepackage[a4paper,margin=1in]{geometry}
\usepackage{amsmath,amssymb,amsthm,microtype}
\newtheorem{theorem}{Theorem}[section]
\theoremstyle{remark}\newtheorem{remark}[theorem]{Remark}
\title{The Critical Shrinking-Annulus Threshold at the Minimal Tail Clock}
\author{Bin Wang}\date{July 2026}
\begin{document}\maketitle
\begin{abstract}
The minimal tail-absorbed clock is shorter than the slope-four clock, but its
noisy tail is only logarithmically small on the target circle.  We determine
how rapidly an auxiliary radius may separate from that circle.  If
$\rho_\sigma=Re^{\eta_\sigma}$ and
$h_\sigma=\lceil a_*\log(1/\sigma)\rceil$, the mass-and-cap noisy-tail scale
is $e^{a_*L_\sigma\eta_\sigma}/L_\sigma$.  For
$\eta_\sigma=c\log L_\sigma/L_\sigma$, the sharp information-class threshold
is $c_*=1/a_*=\log(10/7)$.  A diagonal multiset consisting of repeated
$q$-atoms saturates all three regimes.  The result controls the tail only;
the moving head remains open.
\end{abstract}

\section{Minimal clock and moving radius}
Put
\[
 q=\frac12,\quad R=\frac75,\quad
 a_*=\frac1{\log(1/(qR))}=\frac1{\log(10/7)},\quad
 L_\sigma=\log(1/\sigma),\quad
 h_\sigma=\lceil a_*L_\sigma\rceil.
\]
Let $\rho_\sigma=Re^{\eta_\sigma}$, where $\eta_\sigma\downarrow0$, and
assume $\rho_\sigma<\rho_*$.  For complement traces satisfying
$|\tau_{\sigma,n}|\le4\sigma^{-1}q^n$, define their tail logarithm beyond
$h_\sigma$.

\section{Threshold theorem}
\begin{theorem}\label{thm:threshold}
For $X=H^\infty(\rho_\sigma)$ or $H^2(\rho_\sigma)$, the noisy tail obeys
\[
 \|S_\sigma^{\ge h_\sigma}\|_X
 \le C\frac{\exp(a_*L_\sigma\eta_\sigma)}{L_\sigma}
\]
for all sufficiently small $\sigma$, with $C$ independent of $\sigma$.
The deterministic target tail is smaller:
\[
 \|T^{\ge h_\sigma}\|_X
 \le C'
 \frac{\sigma^{1-a_*\log(q_*/q)}
 \exp(a_*L_\sigma\eta_\sigma)}{L_\sigma}.
\]
If
\[
 \eta_\sigma=c\frac{\log L_\sigma}{L_\sigma},
\]
then the noisy scale is $L_\sigma^{a_*c-1}$.  Thus it vanishes for
$c<c_*=1/a_*$, is of constant order for $c=c_*$, and grows for $c>c_*$.
Here
\[
 c_*=\log(10/7)=0.3566749439387324\ldots.
\]
These three orders are sharp in the mass-and-cap information class.
\end{theorem}
\begin{proof}
Put $x_\sigma=q\rho_\sigma=qR e^{\eta_\sigma}$.  The geometric bounds are
\[
 \|S_\sigma^{\ge h}\|_{H^\infty}
 \le\frac{4\sigma^{-1}x_\sigma^h}{h(1-x_\sigma)},\qquad
 \|S_\sigma^{\ge h}\|_{H^2}
 \le\frac{4\sigma^{-1}x_\sigma^h}
 {h\sqrt{1-x_\sigma^2}}.
\]
Since $(qR)^{a_*L_\sigma}=\sigma$, the first estimate follows; the ceiling
changes only a bounded factor.  The target proof is identical with
$48(q_*\rho_\sigma)^h$ and
\[
 (q_*R)^{a_*L_\sigma}
 =\sigma^{1-a_*\log(q_*/q)}.
\]
The remaining exponent is $0.0531384836\ldots>0$.

For sharpness, take $N_\sigma=\lfloor\sigma^{-1}/q^2\rfloor$ copies of the
atom $q$ on a diagonal operator.  Its squared mass is at most
$\sigma^{-1}$ and its trace powers are $N_\sigma q^n$, asymptotic to the
upper envelope.  The first tail coefficient gives the same lower order in
both norms, while the geometric estimates give the matching upper order.
\end{proof}

\section{Protocol and boundary}
The computation evaluates $c=c_*/2,c_*,3c_*/2$ at $\sigma=10^{-80}$,
records $\eta_\sigma$, $\rho_\sigma$, and the normalized tail scale, and
verifies that every reported radius remains below $\rho_*$.  At the critical
row the normalized model scale equals one by definition; this means
$\Theta(1)$, not convergence of the actual tail to one.

\begin{remark}
The repeated-$q$ diagonal is a saturation model for the available
mass-and-cap information, not the measured noisy spectrum.  The theorem
does not control the moving head and does not prove actual annular
nonconvergence.  Gates A--E remain false/open.
\end{remark}
\end{document}
